\scanpage{4}\mspage{2}
\[
\deleted{2i-1+2j-1=2(i+j-1)=\unsure{2n|2m|n}}
\]

% slanted
\margin{\textquotedblleft{}\uline{Néron}\textquotedblright{}}
\uline{Th. 5}\qquad Soient $X$ variété projective n.s. de dim $n$ \deleted{\illegible{2}}, $Z^i,Z'^{j}\in\mathcal{Z}^{\deleted{\unsure{*|n}}}(X)$\inserted{ de codim $i,j$, avec \deleted{\unsure{\illegible{1}|codim}} $i+j=n$}. Supposons que $z=Z\mathbin{.}Y$ et $z'=Z'\mathbin{.}Y$ soient alg. équiv. à $0$\inserted{ sur $Y$} [$Y$ section hyperplane\inserted{ générique d'un pinceau\inserted{ dont \deleted{toutes} les fibres sont géom. intègres\ldots}}] ; ils sont de codim \deleted{\unsure{$m$ et $m$}}\inserted{$i,j$}, avec \deleted{$\dim Y=2m-1$}\inserted{$i+j=n=\dim Y+1$} donc \deleted{\unsure{ce sont des|définissent des}}\inserted{définissent} des éléments de $P^{\deleted{\unsure{m}}i}(Y),P^{j}(Y)$ et par la forme de Néron-Tate [définie grâce à l'accouplement entre $P^{i}(Y)$ et $P^{j}(Y)$] on \deleted{\unsure{a donc}} définit $\langle z,z'\rangle$. Ceci posé, on a \deleted{[au signe près ??]}
\[
\langle z,z'\rangle=(-1)^m Z\mathbin{.}Z'
\]

% slanted; circled
\margin{\uline{N.B} Sauf si $2i=2j=n$ les deux \unsure{mbs|mbres|membres} sont nuls par Lefschetz \unsure{aux|au|en} \unsure{deux genres|tous genres|sens grossier}}
% A triangular border encloses the note. An arrow points up to Ceci.
\margin{\tikz[baseline=(n.base)]{\node[text width=2cm,align=left,inner sep=1pt] (n) {Doit résulter d'un énoncé plus général sur des fibrations, et admet une \uline{variante} locale};\draw (n.south west) -- (n.north west) -- ([yshift=3mm]n.north east) -- (n.south east) -- cycle;\draw[->] (n.north) -- ++(0,0.3);}}

\deleted{\uline{\unsure{Cor.|Corr.}}} \uline{Th. 6}\qquad Sous les conditions du th 5,

% slanted
\margin{\textquotedblleft{}\uline{Lefschetz}\textquotedblright{}}
Si $2i\leq n-1$ \ $P^i(X)\mathrel{-}P^i(Y)$ est une \uline{isogénie} ; \deleted{\unsure{et}} \uline{si $2i=n$} le \uline{facteur fixe} de $P^i(Y)$ est identique à l'image de $P^i(X)$ ; \deleted{\unsure{[} \deleted{\unsure{et}} C'est égal à $P^i(X)$} \deleted{sauf [Si \deleted{\unsure{pour}} \unsure{$2i=n$} $n=2i$ \unsure{ou} $n$ pair et $i=\frac{n}{2}$]} \deleted{Le noyau de $P^i(X)\mathrel{-}P^i(Y)$ est fini si $2i\leq n-1$} [Si \deleted{$\unsure{2i|2}$} $2i\geq n+1$, $P^i(X)\mathrel{-}P^i(Y)$ est \inserted{surjectif}]
% archive sheet number: 3
